% answers_short.tex —— 1983 年·全国卷（理科） · 简短答案速查
% 由 tools/make_short_answers.py 自 answers.tex 的 \gksolhead/\gkanswer 提取，
% 并与 1983/1983_数学真题_理_答案要点.md 逐题比对后生成，勿手改（改源 MD 后重跑）。

\begin{gknotice}
\textbf{使用说明}：本节只给最终答案便于快速核对，完整推导见解析版。答案为编写者独立推导，\textbf{非官方标准答案}。
\end{gknotice}

\gkpart{题目}
\begin{itemize}[leftmargin=2.6em]
  \item[\textbf{1.}] D
  \item[\textbf{2.}] A
  \item[\textbf{3.}] D
  \item[\textbf{4.}] C
  \item[\textbf{5.}] C
  \item[\textbf{6.}] (1) 两图形交于 $(0,0)$ 与 $(1,-1)$；(2) 表示以 $\left(\dfrac52,0\right)$ 为圆心、$\dfrac52$ 为半径的圆（过极点，直径落在极轴上）
  \item[\textbf{7.}] (1) $\mathrm dy=\mathrm e^{-x}(2\cos 2x-\sin 2x)\,\mathrm dx$；(2) $100$ 种
  \item[\textbf{8.}] $0$
  \item[\textbf{9.}] (1) 见解析版；(2) $k\pi-\dfrac{\pi}{4}\leqslant t\leqslant k\pi+\dfrac{\pi}{4}$（$k\in\mathbf{Z}$，即 $\lvert\tan t\rvert\leqslant1$）
  \item[\textbf{10.}] 见解析版
  \item[\textbf{11.}] $\alpha=\dfrac{\pi}{6}$ 或 $\alpha=\dfrac{5\pi}{6}$
  \item[\textbf{12.}] (1) 见解析版；(2) $\dfrac{b^2}{1+p}$
  \item[\textbf{13.}] (1) 见解析版；(2) 见解析版
\end{itemize}
